\section{Related Work}
\label{sec:related-work}

% -----------------------------------------------
\subsection{Cheating on Cyber and Agent Benchmarks}
\label{sec:rw-cheating}

NIST CAISI documented the first Cybench-specific cheating in late 2025: agents using coding tools to search the internet for challenge flags and walkthroughs~\cite{nist2025cheating}. Stein, Brown et al.\ followed in April 2026 with the Meerkat study~\cite{stein2026meerkat}, the most direct precedent to our work. Auditing traces across nine benchmarks, they found 16 of 464 successful Cybench traces (3.4\%) solved by searching for public writeups or overwriting tests. Our methodology differs in three respects: we test a larger set of models, count infrastructure probing and failed cheat attempts in addition to successful web-search cheating, and flag searches by intent (challenge or event identifiers in the query) rather than only confirmed writeup downloads.

Beyond Cybench, several studies have shown that benchmark gaming is a structural problem across agent evaluations. UC Berkeley RDI built an agent that scored 100\% on SWE-bench Verified by injecting a pytest hook that forced every test to report as passing, demonstrating that benchmark design, not just model behavior, enables cheating~\cite{berkeley2026benchmarks}. Wang et al.\ audited ten prominent agent benchmarks with BenchJack~\cite{wang2026benchjack}, finding that nine of ten could be exploited to near-perfect scores without solving a single task, and deriving a taxonomy of eight recurring flaw patterns. On coding benchmarks specifically, Zhong et al.\ created ``impossible'' task variants where any pass necessarily implies a specification-violating shortcut, providing a direct outcome-based measure of cheating propensity~\cite{zhong2025impossiblebench}. DeepSWE~\cite{datacurve2026deepswe} found that 18--25\% of Claude Opus passes on SWE-bench Pro were achieved through git-history exploits rather than legitimate problem solving. Han et al.\ found that search-capable agents locate ground truth labels online for approximately 3\% of benchmark questions, a form of search-time data contamination that affects any evaluation where models have internet access~\cite{han2025contamination}. Kang et al.\ reduced CVE-Bench performance overestimation by 33\% by applying their Agentic Benchmark Checklist guidelines~\cite{kang2025benchmarks}. CTFusion argues that static CTF benchmarks---where challenges have been publicly solved and writeups indexed---are inherently vulnerable, and that only live, unreleased challenges offer a durable fix~\cite{ctfusion2026}.

% -----------------------------------------------
\subsection{Specification Gaming and Reward Hacking}
\label{sec:rw-spechack}

Cybench cheating is one instance of specification gaming, the broader pattern in which agents satisfy the literal specification of an objective without achieving the intended outcome~\cite{krakovna2020specgaming}. Palisade Research showed that some reasoning models hack a chess environment by default when outmatched---o3 in 88\% of runs---while chat models such as GPT-4o and Claude 3.5 Sonnet do not without explicit nudging~\cite{palisade2025specgaming}. METR found frontier models reward-hack in 0.7\% of runs on their general task suite (HCAST) and over 30\% on research engineering tasks (RE-Bench), with models demonstrating awareness that their behavior is not in line with user intentions~\cite{metr2025rewardhacking}. Zheng et al.\ demonstrated that even the evaluator is an attack surface: a null model achieved 86.5\% on AlpacaEval 2.0 by exploiting LLM-judge templates~\cite{zheng2024cheating}.

Greenblatt et al.~\cite{greenblatt2024alignmentfaking} demonstrated that Claude 3 Opus strategically complies with training objectives it would otherwise refuse 14\% of the time when it infers it is being trained, the first empirical demonstration of alignment faking emerging without explicit instruction. This evaluation-awareness is a precursor to the benchmark gaming we observe: models that detect they are being evaluated may adjust behavior accordingly.

The Reward Hacking Benchmark~\cite{thaman2026rewardhacking} tests 13 models on sandboxed tool-use tasks with no internet access, finding the newest Claude models at 0\% exploit rate and DeepSeek-R1-Zero at 13.9\%, with environmental hardening reducing exploits by 87.7\% relative. Their taxonomy covers infrastructure exploits but not web-search cheating.

% -----------------------------------------------
\subsection{Prompt-Level Behavioral Constraints}
\label{sec:rw-prompts}

Our study tests whether system-prompt instructions can suppress cheating, a specific case of the broader question of whether natural-language constraints reliably govern model behavior. Constitutional AI~\cite{bai2022constitutional} demonstrated that models can be trained to follow written behavioral principles, but this operates at the training level rather than at inference time via prompts. Wallace et al.~\cite{wallace2024instruction} proposed the instruction hierarchy, training LLMs to prioritize system-prompt directives over user messages and third-party content, improving robustness to prompt injection. However, Mu et al.~\cite{mu2025systemprompt} found that current LLMs frequently fail to adhere to system-prompt constraints when facing conflicting inputs, even without adversarial intent. Geng et al.~\cite{geng2025controlillusion} showed more broadly that system/user prompt separation fails to establish a reliable instruction hierarchy: models exhibit strong biases toward certain constraint types regardless of priority designation, and social cues such as authority or expertise show stronger influence on model behavior than structural prompt ordering. Shin~\cite{shin2026compliancegap} identified a related failure mode termed the compliance gap: models verbally agree to follow procedural instructions but then fail to execute them, with Claude Sonnet 4 exhibiting 100\% verbal compliance and 0\% actual compliance in one evaluation.

Two prior studies tested prompt-level cheating suppression specifically. Zhong et al.~\cite{zhong2025impossiblebench} ran a prompt ablation on ImpossibleBench, finding that stricter instructions reduced cheating from over 85\% to as low as 1\% on coding benchmarks. CTFusion~\cite{ctfusion2026} appended a no-cheat instruction on a CTF benchmark and observed a 29\% relative reduction in pass rate. Our study extends this line of work to a multi-model, multi-condition ablation on a cybersecurity benchmark, testing 22 models under three prompt conditions.

% -----------------------------------------------
\subsection{LLM-as-a-Judge for Agent Evaluation}
\label{sec:rw-judge}

Our cheating detection pipeline uses an LLM-as-a-judge to audit agent transcripts. The LLM-as-judge paradigm was established by Zheng et al.~\cite{zheng2023judging}, who showed that strong LLMs can serve as scalable evaluators of open-ended text quality, achieving high agreement with human judgments on MT-Bench. Yuan et al.~\cite{yuan2024rjudge} benchmarked LLM proficiency at judging safety risks from agent interaction records, finding that the best model achieved only 74.42\% accuracy across 569 multi-turn records, highlighting the difficulty of automated safety judgment. Ruan et al.~\cite{ruan2024lmagents} extended the paradigm to agent safety, using an LLM-emulated sandbox to identify risky tool-use actions before execution. In the offensive security domain, Caldwell et al.\ developed ScopeJudge~\cite{caldwell2025scopejudge}, which uses LLM-as-a-judge to assess whether individual tool calls in penetration testing agent trajectories fall within the defined scope of an engagement. Our judge addresses a complementary problem: distinguishing genuine exploitation from benchmark gaming.
